\documentclass[10pt,a4paper]{amsart}
\usepackage[margin=19mm]{geometry}
\usepackage{amsmath,amssymb,mathtools}
\usepackage[scr=rsfs,cal=euler]{mathalfa}
\usepackage{xfrac}
\usepackage[unicode,hidelinks,bookmarks=false]{hyperref}
\newcommand{\ie}{i.e.\ }
\newcommand{\cf}{cf.\ }
\newcommand{\resp}{respectively\ }
\newcommand{\Subsection}{\subsection}
\providecommand{\nobreakdash}{\nobreak\hbox{-}}
\setlength{\emergencystretch}{2em}
\multlinegap0pt
\usepackage{amssymb}
\usepackage{enumerate}
\newtheorem{corollary}{Corollary}
\newtheorem{lemma}{Lemma}
\newtheorem{theorem}{Theorem}
\newcommand{\Fqt}{\mathbb{F}_{q^2}}
\newcommand{\I}{\mathcal{I}}
\def\PP{{\mathbb P}} 
\newcommand{\U}{\mathcal{U}}
\title{Functional codes arising from rank $n$ Hermitian varieties and hypersurfaces in low dimensions}
\author{Subrata Manna}
\date{Source: arXiv:2605.23221v1 (CC BY 4.0)}
\begin{document}
\maketitle

\noindent\emph{Source and licence.} Functional codes arising from rank $n$ Hermitian varieties and hypersurfaces in low dimensions. Version arXiv:2605.23221v1 (CC BY 4.0), distributed by arXiv under the Creative Commons Attribution 4.0 International licence, http://creativecommons.org/licenses/by/4.0/, as recorded in the arXiv licence metadata for this exact version and shown on the abstract page https://arxiv.org/abs/2605.23221v1. Full licence notice: \texttt{licenses/arxiv-2605.23221-CC-BY-4.0.txt}.\par\smallskip

\noindent\textit{Editorial scope.} Counted unit: the article's theorem n (source0.tex lines 291--293). The article's complete original proof of it is delivered with it (source0.tex lines 295--333, one \texttt{proof} environment of 39 lines). No mathematics is written, repaired or re-derived: the statement and the proof are the article's own lines, in the article's own notation, and no retained line is substituted or re-flowed.  Retained uncounted as the source-local context the proof needs: lines 175--177; lines 261--263; lines 270--272.  Nothing was omitted from any retained span, so the package carries no omission record.  No retained line contains a citation, so the wrapper carries no bibliography and no bibliography entry.  No retained line contains a figure environment, an image-inclusion command or drawing code, so no image is delivered and the package declares no asset dependency.  Editorial formatting only, in addition: an AMS \texttt{amsart} preamble that restates the article's own declarations and macros used by the retained lines, namely the environments \texttt{corollary}, \texttt{lemma}, \texttt{theorem} on separate counters, together with the standard packages those lines need.  The retained environments print in this document as Corollary 1, Lemma 1, Theorem 1 in order of appearance, whereas the article prints the same items with the numbering of its own full document.  The complete native source of the article as distributed in the arXiv e-print for this exact version is bundled unchanged as \texttt{sources/201195/source0.tex}.
\begin{corollary}\label{hyp cup}
      A degenerate Hermitian variety $P\U_{n-1}$ of rank $n$ in $\PP^n$ and an $\Fqt$-hyperplane $\Sigma$ with $P\notin \Sigma$ intersect in a non-degenerate Hermitian variety $\U_{n-1}$ in $\Sigma\cong \PP^{n-1}$.
\end{corollary}

\begin{lemma}\label{outside hyperplane}
 Let $n\geq 3$. If $\Sigma$ is an $\Fqt$-hyperplane in $\PP^n$  such that $P \not\in \Sigma$ and $\Sigma \subseteq  V(F)$, then $|P\U_{n-1} \cap V(F)\cap \Sigma^C(\Fqt)| \le (d-1)(q+1)q^{2n-4}$, where $\Sigma^C = \PP^n \setminus \Sigma$. 
\end{lemma}

\begin{lemma}\label{missing P}
  Let $n\geq 3$. If $P\notin V(F)$, then $$|P\U_{n-1}\cap V(F)(\Fqt)|\leq d|\U_{n-1}(\Fqt)|< |\U_{n-1}(\Fqt)|+(d-1)(q+1)q^{2n-4}.$$   
\end{lemma}

\begin{theorem}\label{rank n}
  Let $n\geq 3$. We have $$|P\U_{n-1}\cap V(F)(\Fqt)|\leq \max \left\{|\U_{n-1}(\Fqt)|+(d-1)(q+1)q^{2n-4}, \ 1+q^2 M_{n-1}(d) \right\},$$ where $M_{n-1}(d):= \max_{V(G)}|\U_{n-1}\cap V(G)(\Fqt)|$ and the maximum is taken over all the $\Fqt$-hypersurfaces $V(G)$ of degree $d\leq q$ in $\PP^{n-1}$.   
\end{theorem}

\begin{proof}
The proof is divided into two cases, according as $P\in V(F)$ or $P\notin V(F)$.

    \textbf{Case 1: $P\notin V(F)$.} It follows from Lemma \ref{missing P} that $$|P\U_{n-1}\cap V(F)(\Fqt)|< |\U_{n-1}(\Fqt)|+(d-1)(q+1)q^{2n-4}.$$

    \textbf{Case 2: $P\in V(F)$.} We distinguish two subcases according to whether the hypersurface $V(F)$ contains an $\Fqt$-hyperplane that does not pass through the point $P$.


    \textit{Subcase 2a.} \textit{$V(F)$ contains an $\Fqt$-hyperplane not passing through $P$.} Suppose $V(F)$ contains an $\Fqt$-hyperplane $\Sigma$ such that $P\notin \Sigma$. It then follows from Corollary \ref{hyp cup} that the hyperplane $\Sigma$ intersects the singular Hermitian variety $P\U_{n-1}$ at a non-degenerate Hermitian variety $\U_{n-1}$. Since $\Sigma\subset V(F)$, we have $P\U_{n-1}\cap V(F)\cap \Sigma =\U_{n-1}$. Consequently, 
    \begin{align*}
        |P\U_{n-1}\cap V(F)(\Fqt)| &= |P\U_{n-1}\cap V(F)\cap \Sigma(\Fqt)|+|P\U_{n-1}\cap V(F)\cap \Sigma^C(\Fqt)|\\
        & \leq |\U_{n-1}(\Fqt)|+(d-1)(q+1)q^{2n-4},
        \end{align*}
        where the last inequality follows from Lemma \ref{outside hyperplane}.

        \textit{Subcase 2b.} \textit{$V(F)$ does not contain any $\Fqt$-hyperplane not passing through $P$.} We now define the sets 
        \begin{align*}
            \Gamma_1 &:=P\U_{n-1}\cap V(F)(\Fqt)\setminus \left\{P\right\} \ \text{and}\\
            \Gamma_2 &:= \left\{\Sigma: \Sigma \ \text{is an $\Fqt$-hyperplane},\ P\notin \Sigma \right\}.
        \end{align*}
Since there are $|\PP^{n-1}(\Fqt)|$ many $\Fqt$-hyperplanes through the point $P$ in $\PP^n$, it follows that $|\Gamma_2|=q^{2n}$. Consider the incidence set $$\I:=\left\{\left(Q,\Sigma\right)\in \Gamma_1\times \Gamma_2 \ | \ Q\in \Sigma  \right\}.$$ For a chosen point $Q\in \Gamma_1$ we see that 
\begin{align}
\#\left\{ \Sigma\in \Gamma_2: Q\in \Sigma \right\} &= \#\left\{ \Sigma\ |\  Q\in \Sigma \right\}-\#\left\{ \Sigma\ |\ \text{the line}\  PQ\subset \Sigma   \right\}\notag\\
         &=|\PP^{n-1}\left(\Fqt\right)|- |\PP^{n-2}\left(\Fqt\right)|\notag \\
         &=q^{2n-2}\label{count hyp}.
\end{align}
We will count $|\I|$ in two ways. First, we observe that, 
\begin{equation}\label{first count}
  |\I|=\sum_{Q\in \Gamma_1}\#\left\{ \Sigma\in\Gamma_2: Q\in \Sigma \right\}=|\Gamma_1|q^{2n-2},  
\end{equation}
where the last equality follows from \eqref{count hyp}. On the other hand, since $P\notin \Sigma$ for any hyperplane in $\Gamma_2$, it follows that $\Sigma\cap P\U_{n-1}$ is a non-degenerate Hermitian variety $\U_{n-1}$ in $\PP^{n-1}$. Moreover, since by assumption $V(F)$ does not contain any hyperplane $\Sigma$ of $\Gamma_2$, the intersection $V(F)\cap \Sigma$ is an $\Fqt$-hypersurface of degree $d$ in $\Sigma\cong\PP^{n-1}$. Consequently, we count that 
\begin{align}
    |\I|&=\sum_{\Sigma\in \Gamma_2}\#\{ Q\in \Gamma_1: Q\in \Sigma \}\notag\\
        &= \sum_{\Sigma\in \Gamma_2} |\Sigma\cap P\U_{n-1}\cap V(F)\left(\Fqt\right)| \notag\\
        &=\sum_{\Sigma\in \Gamma_2}|\U_{n-1}\left(\Fqt\right)\cap \left(\Sigma\cap V(F)  \right)|\notag\\
        &\leq |\Gamma_2| M_{n-1}(d)=q^{2n}M_{n-1}(d).\label{second count}
\end{align}
Thus it follows from comparing \eqref{first count} and \eqref{second count} that $|\Gamma_1|\leq q^2 M_{n-1}(d)$, which implies that $$|P\U_{n-1}\cap V(F)(\Fqt)|\leq 1+q^2M_{n-1}(d).$$ Now the theorem follows from the two cases.
\end{proof}

\end{document}
